\appendix

\section{Outlier and Sensitivity Cases}
\label{app:outliers}

Table~\ref{tab:outliers} lists cases that require cautious interpretation. These cases were not removed from the analysis because they are informative for understanding dataset- and model-dependent split effects.

\begin{table}[H]
\centering
\caption{Outlier and sensitivity cases identified during the split-diagnostic audit.}
\label{tab:outliers}
\resizebox{\textwidth}{!}{%
\begin{tabular}{llllrrl}
\toprule
Dataset & Task & Model & Metric & Ordinary gap & Balanced gap & Interpretation flag \\
\midrule
ClinTox & Classification & LogisticRegression & ROC-AUC & 0.119 & 0.054 & Class-imbalance-sensitive dataset \\
ClinTox & Classification & RandomForest & ROC-AUC & 0.094 & 0.046 & Class-imbalance-sensitive dataset \\
ClinTox & Classification & XGBoost & ROC-AUC & 0.033 & 0.057 & Class-imbalance-sensitive dataset \\
ESOL & Regression & RandomForest & RMSE & 1.038 & 0.340 & Large gap change after balancing \\
ESOL & Regression & Ridge & RMSE & -0.082 & 0.793 & Linear model behaves differently from tree models \\
ESOL & Regression & XGBoost & RMSE & 0.974 & 0.207 & Large gap change after balancing \\
FreeSolv & Regression & RandomForest & RMSE & 0.737 & -0.323 & Negative balanced gap; small-dataset constraint \\
FreeSolv & Regression & XGBoost & RMSE & 1.029 & 0.069 & Large gap change after balancing \\
\bottomrule
\end{tabular}%
}
\end{table}

\section{Model Ranking Stability}
\label{app:rankings}

Split choice can change not only absolute performance estimates but also model rankings. This is important because benchmark papers often use model rankings to support claims about method superiority. In ESOL, for example, XGBoost and Random Forest rank highest under random and balanced scaffold splits, whereas Ridge ranks highest under the ordinary scaffold split. In HIV, Random Forest ranks highest under random and balanced scaffold splits, while XGBoost ranks highest under the ordinary scaffold split. These examples show that split diagnostics should accompany model comparisons.

\section{Reproducibility Notes}
\label{app:reproducibility}

The full experimental workflow is implemented in the repository scripts. The main analysis sequence is:

\begin{verbatim}
python paper1_leakage_benchmark/scripts/01_prepare.py
python paper1_leakage_benchmark/scripts/02_featurize_and_split.py
python paper1_leakage_benchmark/scripts/train_many.py
python paper1_leakage_benchmark/scripts/train_balanced_scaffold.py
python paper1_leakage_benchmark/scripts/compare_splits.py
python paper1_leakage_benchmark/scripts/scaffold_stats.py
python paper1_leakage_benchmark/scripts/balanced_scaffold_stats.py
python paper1_leakage_benchmark/scripts/similarity_audit.py
python paper1_leakage_benchmark/scripts/outlier_inspection.py
python paper1_leakage_benchmark/scripts/make_paper_figures.py
python paper1_leakage_benchmark/scripts/make_manuscript_tables.py
\end{verbatim}

All models use the same RDKit Morgan fingerprint representation so that the analysis focuses on the effect of split protocols rather than feature-engineering differences.
